\section{Operaciones básicas e indexación}

\subsection{Series}

\subsubsection{Indexación}

Si usted estuviera interesado en buscar un dato específico de \textit{Pandas} puede utilizar la indexación o slicing de la misma manera que se utilizaba en \textit{NumPy}.

Cuando presentamos el tema de \textit{Series}, utilizamos la indexación, indexar consta de la siguiente estructura:
\begin{verbatim}
    serie[índice]
\end{verbatim}
En código:
\begin{pycode}{Indexación de series en pandas}
import pandas as pd

# Definimos las variables
indice = ['Enero', 'Febrero', 'Marzo']
datos = [200, 100, 135]

# Creamos la serie de Pandas
ingresos_mensuales = pd.Series(datos, index=indice) 

# Indexamos
print('Ingresos del mes de marzo', ingresos_mensuales['Marzo']) # Retorna: 'Ingresos del mes de marzo 135'
# ingresos_mensuales es la serie, [] la indexacion y 'Marzo' el indice
\end{pycode}

\paragraph{Slicing}
Así es, al igual que en \textit{NumPy} el \textit{Slicing} tambien existe en pandas, para realizar slicing, se realiza el mismo proceso que en \textit{NumPy}, 
\begin{verbatim}
    Serie[Inicio:Fin:paso]
\end{verbatim}

\begin{nota}
    \textit{¡Alerta de diferencia!:} En \textit{NumPy} y en las listas de \textit{Python}, el límite superior (\textit{Fin}) nunca se incluye. Pero en \textit{Pandas}, cuando hacés \textit{slicing} con \textbf{etiquetas} de texto, el límite final SÍ se incluye.
\end{nota}

Veamoslo en un ejemplo de código:
\begin{pycode}{Slicing en series}
import pandas as pd

# Definimos las variables
indice = ['Enero', 'Febrero', 'Marzo']
datos = [200, 100, 135]

# Creamos la serie de Pandas
ingresos_mensuales = pd.Series(datos, index=indice) 

# Realizamos Slicing
print(ingresos_mensuales['Enero':'Febrero']) # Retorna:
# Enero      200
# Febrero    100
# dtype: int64
\end{pycode}

\subsubsection{Operaciones matemáticas básicas}

En una serie puede realizar operaciones matemáticas básicas de:
\begin{itemize}
    \item Aritméticas: suma, multiplicación\dots
    \item Comparaciones: de mayor o menor (retorna True o False)
    \item Estadísticas: de media, valor maximo, etc.
\end{itemize}

Veamos algun ejemplo básico
\begin{pycode}{Operaciones matemáticas}
import pandas as pd

# Definimos las variables
indice = ['Enero', 'Febrero', 'Marzo']
datos = [200, 100, 135]

# Creamos la serie de Pandas
ingresos_mensuales = pd.Series(datos, index=indice) 

# Operación aritmética
print(ingresos_mensuales['Enero'] + 10) # Retorna: 210

# Operación de comparación
print(ingresos_mensuales > 120) # Retorna:
# Enero       True
# Febrero    False
# Marzo       True

# Operación Estadística
print(ingresos_mensuales.mean()) # Retorna: 145
\end{pycode}

\begin{nota}
    Si usted decide no aclarar el indice en la operación, la operación se va a aplicar a todo el data frame. Esto se debe a que \textit{Pandas} utiliza las operaciones vectorizadas de \textit{NumPy}
\end{nota}




\subsection{Data Frames}

\subsubsection{Indexación}

Hay diferentes formas de indexar un \textit{Data Frame}. 

\paragraph{Selección Directa (Columnas)}

Es la forma más rápida de ``recortar'' una tabla verticalmente.
\begin{itemize}
\item \textit{Un solo corchete} \texttt{df['header']}: Extrae una Serie. Es útil cuando solo quieres operar sobre una variable (ej: calcular el promedio de impuestos).
\item \textit{Doble corchete} \texttt{df[['h1', 'h2']]}: Extrae un nuevo DataFrame. Piénselo como "pedir una lista de columnas". 
\end{itemize}

\paragraph{Indexación por \texttt{.loc} y \texttt{.iloc}}

El atributo \texttt{.loc} funciona con los nombres de las cabezeras (headers) de las columnas y en las etiquetas de las filas (indices).

Si usted quisiera realizar un slicing con \texttt{.loc}, el límite final si se incluye. Es decir si usted realiza ``\texttt{df.loc['A':'C']}'' obtendria la serie A, B y C. 

Su estructura consiste en:
\begin{verbatim}
    dataframe.loc[NombreFila, NombreColumna]
\end{verbatim}

El atributo \texttt{.iloc} es posicional, es decir que no ve las etiquetas, sino el orden en que los datos fueron cargados en la memoria. A diferencia de \texttt{.loc} este NO incluye el límite final. 

\textit{.iloc} permite utilizar numeros negativos, por ejemplo utilizar \texttt{df.iloc[-5:]} retornaria las ultimas cinco filas.

Su estructura consiste en:
\begin{verbatim}
    dataframe.loc[NúmeroFila, NúmeroColumna]
\end{verbatim}


\paragraph{.head y .tail}
Si usted quisera ver las primeras filas, puede utilizar \texttt{.head()}. Dentro del parentesis puede colocar hasta donde que fila quiere que se muestre. (por defecto es 5). \texttt{.tail()} funciona igual que \texttt{.head()} pero retorna las ultimas filas.

Veamos todas las formas de indexar y de realizar slicing de un data frame en código:
\begin{pycode}{Indexación y Slicing de un data frame}
import pandas as pd

# Creamos un diccionario:
mi_mercado = {
    'productos': ['bananas', 'peras', 'manzanas'], 
    'ingresos_mensuales': [100,60,70],
    'impuestos': [0.1,0.05,0.07]
}

# Creamos el Data Frame
data_frame = pd.DataFrame(mi_mercado)


print(data_frame['productos']) # Retorna
# 0     bananas
# 1       peras
# 2    manzanas
print()

print(data_frame[['productos', 'impuestos']]) # Retorna
#   productos  impuestos
# 0   bananas       0.10
# 1     peras       0.05
# 2  manzanas       0.07
print()

print(data_frame.loc[0:2]) # Retorna
#   productos  ingresos_mensuales  impuestos
# 0   bananas                 100       0.10
# 1     peras                  60       0.05
# 2  manzanas                  70       0.07

# Note que incluye el indice 2
print()

print(data_frame.iloc[0:2]) # Retorna
#   productos  ingresos_mensuales  impuestos
# 0   bananas                 100       0.10
# 1     peras                  60       0.05
print()

print(data_frame.head(1)) # Retorna
#   productos  ingresos_mensuales  impuestos
# 0   bananas                 100        0.1
\end{pycode}

\begin{ejemplo}
    Si usted quisiera indexar por filas y columnas podria utilizar: \texttt{data\_frame['productos'].loc[1]} que le retornaria 'peras'
\end{ejemplo}



\subsubsection{Operaciones Matemáticas}

En \textit{pandas} las operaciones vectorizadas se pueden realizar entre columnas. Hay diferentes tipos de operaciones matemáticas que se puede realizar, pero veamoslas en código.

\paragraph{Operaciones aritméticas}
\begin{pycode}{Operaciones aritméticas en data frames}
import pandas as pd

# Creamos el data frame
df = pd.DataFrame({
'Producto': ['Soja', 'Maíz', 'Trigo'],
'Toneladas': [100, 200, 150],
'Precio_Unitario': [400, 250, 300]
})

# Multiplicación entre columnas
df['Ingreso_Bruto'] = df['Toneladas'] * df['Precio_Unitario'] # Nos retorna el Ingreso Total

# Operación con un escalar (Aplicar un impuesto del 10%)
df['Precio_con_IVA'] = df['Precio_Unitario'] * 1.10

# Suma y Resta (Calcular stock neto tras una perdida de 5 tons)
df['Stock_Neto'] = df['Toneladas'] - 5
\end{pycode}

\paragraph{Operaciones de comparación}
\begin{pycode}{Comparacion en Data Frames}
import pandas as pd

# Creamos el data frame
df = pd.DataFrame({
'Producto': ['Soja', 'Maíz', 'Trigo'],
'Toneladas': [100, 200, 150],
'Precio_Unitario': [400, 250, 300]
})

# Multiplicación entre columnas
df['Ingreso_Bruto'] = df['Toneladas'] * df['Precio_Unitario'] # Nos retorna el Ingreso Total

# Operación comparativa para ver ingresos mayores de 40000
mascara = df['Ingreso_Bruto'] > 40000 # Retorna: [False, True, True]

#Usar la comparacion para filtrar el DataFrame original
df_vip = df[df['Ingreso_Bruto'] > 40000] # Retorna
#   Producto  Toneladas  Precio_Unitario  Ingreso_Bruto
# 1     Maíz        200              250          50000
# 2    Trigo        150              300          45000
\end{pycode}


\paragraph{Operaciones estadísticas}

\begin{pycode}{Estadistica en Pandas}
import pandas as pd

# Creamos el data frame
df = pd.DataFrame({
'Producto': ['Soja', 'Maíz', 'Trigo'],
'Toneladas': [100, 200, 150],
'Precio_Unitario': [400, 250, 300]
})

# Operación suma
total_tons = df['Toneladas'].sum()

# Operación para obtener la media
promedio_precio = df['Precio_Unitario'].mean()

# Devuelve todas las estadísticas
resumen = df.describe()
\end{pycode}


\begin{nota}
    \textit{¡Atención!}: Aunque puede usar ``+'' o ``*'', Pandas tiene métodos como \texttt{.add()} o \texttt{.mul().}. Esto porque permiten manejar valores faltantes \texttt{(NaN)}. 
    
    Si una celda está vacía, el operador ``+'' dará error o NaN, pero con \texttt{.add(fill\_value=0)} puede decirle a \textit{Pandas} que trate el vacío como un cero para no romper el cálculo.
\end{nota}